\chapter{문제 정의와 제안 방법}\label{chap:method}

이 장에서는 먼저 Neuron Chunking 문제의 최적해가 단일 연속 구간이 됨을 보이고, 문제를 정확히 $R$개 행을 선택하는 형태로 바꾼다. 이어서 SSD의 처리량 포화를 반영한 two-line 지연 모델을 도입하고, 이 모델 아래에서 최적 mask를 구하는 동적계획법과 계산이 가벼운 선택 방법인 Tiles, 그리고 선택한 mask가 최적값에 얼마나 가까운지에 대한 보장을 차례로 제시한다.

\section{Neuron Chunking 문제의 최적해}\label{sec:single-chunk}

\ref{sec:neuron-chunking}절의 Neuron Chunking 문제를 다시 생각해보자.
\[
\max_{\substack{M\in\{0,1\}^N\\1\le\|M\|_1\le R}}
\frac{I(M)}{\widehat L_{\mathrm{table}}(M)}.
\]
$1\le R\le N$이고, 모든 길이 $1\le\ell\le R$에 대해 $T[\ell]>0$이라 하자.

행 예산 $1\le\|M\|_1\le R$ 아래에서 임의의 최적 mask $M^\star$를 생각하자. $\mathcal C(M^\star)$의 구간들은 선택된 행들을 분할하므로
\[
\begin{aligned}
\frac{I(M^\star)}{\widehat L_{\mathrm{table}}(M^\star)}
&=
\sum_{C\in\mathcal C(M^\star)}
\underbrace{
\frac{T[|C|]}{\widehat L_{\mathrm{table}}(M^\star)}
}_{w_C}
\underbrace{
\frac{\sum_{i\in C}v_i}{T[|C|]}
}_{U(C)}\\
&=
\sum_{C\in\mathcal C(M^\star)}w_CU(C)
\le
\max_{C\in\mathcal C(M^\star)}U(C).
\end{aligned}
\]
$w_C>0$이고 $\sum_Cw_C=1$이므로, 전체 비율은 각 구간의 비율을 latency로 가중한 평균이다. 따라서 가장 큰 구간 비율을 넘을 수 없다.

$U(C)$가 가장 큰 구간 $C^\star\in\mathcal C(M^\star)$만 남긴 mask를 $\widetilde M$이라 하자.

\[1\le\|\widetilde M\|_1=|C^\star|\le\|M^\star\|_1\le R\]

이므로 $\widetilde M$도 행 예산을 만족한다. 앞선 부등식에 의해 $\widetilde M$의 비율은 $M^\star$보다 작지 않고, $M^\star$의 최적성에 의해 더 클 수도 없다. 따라서 $\widetilde M$도 최적해다.

따라서 단일 연속 구간으로 이루어진 최적해가 존재하며, 원래 문제는 길이가 $R$ 이하인 모든 연속 구간 중 $U(C)$를 최대화하는 구간을 찾는 문제가 된다.
\[
\max_{\substack{M\in\{0,1\}^N\\1\le\|M\|_1\le R}}
\frac{I(M)}{\widehat L_{\mathrm{table}}(M)}
=
\max_{\substack{1\le\ell\le R\\0\le j\le N-\ell}}
U(C_{j,\ell}).
\]
오른쪽의 비교 대상은 후보 집합 $\mathcal G$에 한정되지 않는 모든 연속 구간이다.

이러한 현상은 목표함수가 정확히 $R$개의 행을 선택하는 것이 아니라, 행 수의 상한에 제약을 거는 것으로 느슨하게 정의되어 있기 때문에 발생한다. 따라서 본 연구에서는 정확히 $R$개의 행을 선택하는 더 강한 문제에 대해 논의를 이어가고자 한다.
\[
\boxed{
\max_{\substack{M\in\{0,1\}^N\\\|M\|_1=R}}
\frac{I(M)}{\widehat L(M)}.
}
\]

\section{Two-line 지연 모델}\label{sec:two-line}

\subsection{연속 구간의 읽기 비용}\label{sec:two-line-chunk}

SSD의 처리량은 요청 크기가 증가함에 따라 증가하다가 일정 수준에서 포화되는 경향을 보인다. 포화 구간의 행당 읽기 비용을 $\alpha>0$라 두면, 길이가 $\ell$행인 I/O 요청의 지연시간을 다음과 같이 근사할 수 있다.
\[
T[\ell]=\alpha\ell,
\qquad \ell\gg1.
\]
작은 요청에서의 처리량 감소를 함께 반영하기 위해, 본 연구에서는 측정표 $T[\ell]$를 다음의 two-line 모델로 근사한다.

길이가 $\ell$행인 구간의 비용을 $T_{\mathrm{2L}}(\ell)$이라 하자. 행 단위로 환산한 전환점을 $\ell_0>0$, 추가 비용의 계수를 $0\le\lambda\le\alpha$라 두면
\[
\boxed{
\begin{aligned}
T_{\mathrm{2L}}(\ell)
&=\alpha\ell+\lambda(\ell_0-\ell)_+\\
&=
\begin{cases}
(\alpha-\lambda)\ell+\lambda\ell_0,
&0<\ell<\ell_0,\\
\alpha\ell,
&\ell\ge\ell_0.
\end{cases}
\end{aligned}
}
\]
여기서 $(x)_+=\max(x,0)$이며, 읽을 행이 없는 경우는 별도로 $T_{\mathrm{2L}}(0)=0$이라 둔다.

이때 두 직선은 $\ell=\ell_0$에서 연속이며, 행 단위 처리량은
\[
\frac{\ell}{T_{\mathrm{2L}}(\ell)}
=
\frac{1}{\alpha+\lambda(\ell_0/\ell-1)_+}
\le\frac1\alpha
\]
이다. $\lambda>0$이면 전환점 이전에는 길이가 늘수록 처리량이 증가하고, 전환점 이후에는 $1/\alpha$로 일정해진다.

모델의 계수 $\alpha,\lambda$와 전환점 $\ell_0$는 SSD 읽기 비용 측정값에 적합하여 결정한다.

\subsection{고정된 행 예산에서의 전체 읽기 비용}\label{sec:two-line-total}

이제 이 모델을 이용해 정확히 $R$개 행을 선택한 mask $M$의 전체 읽기 비용을 살펴보자. 선택한 행 수가 같더라도 연속 구간의 구성에 따라 비용이 달라지며, two-line 모델에서는 이 차이를 짧은 구간의 추가 비용으로 표현할 수 있다.

선행연구의 구간별 비용을 합산하는 구조에 two-line 모델을 적용하자. $\sum_{C\in\mathcal C(M)}|C|=R$이므로 전체 비용은 다음과 같이 분해된다.
\[
\begin{aligned}
\widehat L_{\mathrm{2L}}(M)
&=\sum_{C\in\mathcal C(M)}T_{\mathrm{2L}}(|C|)\\
&=\alpha R
+\lambda\sum_{C\in\mathcal C(M)}(\ell_0-|C|)_+.
\end{aligned}
\]
따라서 같은 $R$에서 모델 비용의 차이는 전환점보다 짧은 구간들의 길이 부족분으로 결정된다. 각 구간은 길이가 $\ell_0$에 도달할 때까지 추가 비용을 줄일 수 있으며, 그 이후에는 추가 비용이 없다.

이 분해를 이용하면 고정된 행 예산에서의 비용 하한도 얻을 수 있다. 모든 구간의 길이가 $R$ 이하이므로 $\sum_C(\ell_0-|C|)_+\ge(\ell_0-R)_+$이고, 따라서
\[
\widehat L_{\mathrm{2L}}(M)
\ge
\alpha R+\lambda(\ell_0-R)_+
=T_{\mathrm{2L}}(R).
\]
길이 $R$의 단일 구간은 이 하한을 달성한다. $\lambda>0$일 때 $R<\ell_0$이면 단일 구간만 하한을 달성하지만, $R\ge\ell_0$이면 모든 구간의 길이가 $\ell_0$ 이상인 mask가 하한을 달성한다. 즉, 여러 구간으로 나뉘어 있더라도 각 구간이 충분히 길면 모델 비용은 동일한 $\alpha R$이다.

한편 행의 배치를 바꾸면 보존 중요도 $I(M)$도 달라진다. 따라서 다음 절에서는 중요도와 짧은 구간의 추가 비용을 함께 고려하는 다음 문제의 최적해를 구한다.
\[
\max_{\substack{M\in\{0,1\}^N\\\|M\|_1=R}}
\frac{I(M)}
{\displaystyle
\alpha R+
\lambda\sum_{C\in\mathcal C(M)}(\ell_0-|C|)_+
}.
\]

\section{동적계획법을 이용한 최적 mask 탐색}\label{sec:dp}

앞 절에서 정의한 two-line 비용 모델 아래에서 정확히 $R$개 행을 선택하는 문제를 생각하자. 표기를 줄이기 위해 $T(\ell)=T_{\mathrm{2L}}(\ell)$, $L(M)=\widehat L_{\mathrm{2L}}(M)$이라 쓰고, 최적 비율을
\[
q^\star=
\max_{\substack{M\in\{0,1\}^N\\\|M\|_1=R}}
\frac{I(M)}{L(M)}
\]
이라 하자. 여기서 $1\le R\le N$, $v_i\ge0$이며, 앞 절의 계수 조건에 의해 모든 가능한 mask의 비용은 $L(M)\ge T(R)>0$이다.

비율의 추정값 $q$를 고정하면, 다음의 합 형태 목적함수를 동적계획법(DP) \cite{bellman1957}으로 최대화할 수 있다.
\[
F_R(q)=
\max_{\substack{M\in\{0,1\}^N\\\|M\|_1=R}}
\{I(M)-qL(M)\}.
\]
먼저 고정된 $q$에서 이 문제를 푸는 방법을 설명한 뒤, DP 결과를 이용해 $q^\star$를 찾는 과정을 설명한다.

\subsection{고정된 \texorpdfstring{$q$}{q}에서의 DP}\label{sec:dp-fixed-q}

행의 인덱스는 $0,\ldots,N-1$로 두고, 중요도의 누적합을 $P_j=\sum_{i=0}^{j-1}v_i$, $P_0=0$으로 정의하자. $D_{z,r}$는 앞의 $z+r$개 행에서 $z$개를 제외하고 $r$개를 선택했을 때의 $I-qL$의 최댓값이다. 필요한 상태의 범위는 $0\le z\le N-R$, $0\le r\le R$이다.

$z\ge1$인 mask에서 마지막으로 선택하지 않은 행을 기준으로 나누자. 그 행 앞에서 선택한 행 수가 $t$라면, 앞부분에는 제외한 행이 $z-1$개, 선택한 행이 $t$개 있다. 마지막으로 제외한 행의 인덱스는 $z+t-1$이고, 그 뒤의 $[z+t,z+r)$는 길이가 $r-t$인 연속 선택 구간이다.

제외한 행이 앞부분과 마지막 구간을 분리하므로 두 부분의 비용을 합산할 수 있다. 마지막 구간의 중요도는 $P_{z+r}-P_{z+t}$이므로
\[
\boxed{
D_{z,r}
=\max_{0\le t\le r}
\left\{
D_{z-1,t}
+P_{z+r}-P_{z+t}
-qT(r-t)
\right\}.
}
\]
$t=r$이면 mask가 선택하지 않은 행으로 끝나며, 마지막 구간은 비어 있다. 이 경우 $T(0)=0$을 사용한다. 모든 mask는 이러한 분해를 가지므로, 모든 $t$를 비교하는 위 점화식은 최적값을 구한다.

제외한 행이 없으면 앞의 $r$개 행을 모두 선택해야 하므로 초기값과 최종값은
\[
D_{0,r}=P_r-qT(r),
\qquad
F_R(q)=D_{N-R,R}
\]
이다. 최적 mask는 $(z,r)=(N-R,R)$에서 역추적한다. 각 상태에서 최댓값을 만드는 $t$를 찾아 $[z+t,z+r)$를 선택하고, $(z,r)\leftarrow(z-1,t)$로 이동한다. $z=0$에 도달하면 남은 앞의 $r$개 행을 모두 선택한다.

\subsection{Two-line 구조를 이용한 DP 계산}\label{sec:dp-two-line}

점화식에서 각 상태마다 모든 $t$를 따로 비교하면 계산량이 커진다. Two-line 모델에서는 마지막 구간을 빈 구간, 짧은 구간, 긴 구간으로 나누어 최댓값 계산을 공유할 수 있다.

정수 길이를 구분하기 위한 경계를
\[
h=\min(\lfloor\ell_0\rfloor,R)
\]
로 두자. $1\le r-t\le h$인 경우에는 짧은 구간의 비용식을, $r-t>h$인 경우에는 긴 구간의 비용식을 사용한다. $\ell_0$가 정수이면 경계에서 두 비용식이 같으므로 그 길이를 짧은 구간 쪽에 포함해도 된다. 비용 계산에 쓰는 전환점 $\ell_0$는 반올림하지 않는다.

점화식에서 $t$에 의존하는 항을
\[
\begin{aligned}
A_{z,t}&=D_{z-1,t}-P_{z+t}+q(\alpha-\lambda)t,\\
B_{z,t}&=D_{z-1,t}-P_{z+t}+q\alpha t
\end{aligned}
\]
로 모으면
\[
D_{z,r}=
\max\left\{
\begin{aligned}
&D_{z-1,r},\\
&P_{z+r}-q\bigl((\alpha-\lambda)r+\lambda\ell_0\bigr)
+\max_{\max(0,r-h)\le t<r}A_{z,t},\\
&P_{z+r}-q\alpha r
+\max_{0\le t<r-h}B_{z,t}
\end{aligned}
\right\}.
\]
세 항은 각각 마지막 구간이 비어 있는 경우, 짧은 경우, 긴 경우에 해당한다. 후보가 없는 범위의 최댓값은 $-\infty$로 둔다.

$r$이 증가할 때 $A_{z,t}$를 조회하는 범위는 폭이 $h$인 이동 구간이며, $B_{z,t}$를 조회하는 범위는 오른쪽 끝이 늘어나는 접두 구간이다. 따라서 이동 구간 최댓값과 누적 최댓값을 이용하면, 고정된 $z$에서 모든 $r$의 값을 $O(R+1)$ 시간에 계산할 수 있다.

DP 표의 크기는 $(N-R+1)(R+1)$이다. 역추적에서도 각 $z$마다 최대 $R+1$개의 후보를 비교하므로, 고정된 $q$에서의 DP 계산과 mask 복원에 필요한 시간과 메모리는 모두
\[
O\bigl((N-R+1)(R+1)\bigr)
\]
이다.

\subsection{최적 비율 \texorpdfstring{$q$}{q}의 탐색}\label{sec:dinkelbach}

가능한 mask $M$에 대해
\[
I(M)-qL(M)
=L(M)\left(\frac{I(M)}{L(M)}-q\right)
\]
이고 $L(M)>0$이므로 다음이 성립한다.
\[
F_R(q)
\begin{cases}
>0,&q<q^\star,\\
=0,&q=q^\star,\\
<0,&q>q^\star.
\end{cases}
\]
최적 비율은 $F_R(q)=0$의 해이다. 이 해를 구하기 위해 Dinkelbach 방법 \cite{dinkelbach1967}을 사용한다.

길이 $R$의 연속 구간 중 중요도가 가장 큰 mask를 $M_C$, 중요도가 가장 높은 $R$개 행을 선택한 mask를 $M_T$라 하자.

$M_C$의 비용은 $T(R)$이고, 모든 가능한 mask는 $I(M)\le I(M_T)$, $T(R)\le L(M)$을 만족하므로 초기 하한과 상한을
\[
q_0=\frac{I(M_C)}{T(R)}
\le q^\star\le
u_0=\frac{I(M_T)}{T(R)}
\]
로 둔다.

$k$번째 반복에서 $q=q_k$로 DP를 풀고, 복원한 mask를 $M_k$라 하자. $F_R(q_k)=I(M_k)-q_kL(M_k)$을 이용하여 다음 비율을
\[
\boxed{
\begin{aligned}
q_{k+1}
&=\frac{I(M_k)}{L(M_k)}\\
&=\frac{q_kL(M_k)+F_R(q_k)}{L(M_k)}\\
&=q_k+\frac{F_R(q_k)}{L(M_k)}
\end{aligned}
}
\]
로 갱신한다.

$q_k\le q^\star$이면 $F_R(q_k)\ge0$이고, $M_k$도 정확히 $R$개 행을 선택하므로 $q_k\le q_{k+1}\le q^\star$이다. $F_R(q_k)>0$인 동안 비율은 증가하고, $F_R(q_k)=0$이면 $M_k$는 원래 비율 문제의 최적해이다. 가능한 mask의 수는 유한하므로, 정확한 DP 풀이와 산술 연산 아래에서는 유한 번의 반복으로 최적값에 도달한다.

\subsection{잔차를 이용한 상한의 갱신}\label{sec:upper-bound}

임의의 최적 mask $M^\star$에 대해 $I(M^\star)=q^\star L(M^\star)$이다.

또한 $q^\star-q_k\ge0$이고 $T(R)\le L(M^\star)$이다.

$F_R(q_k)$의 정의에 의해 $I(M^\star)-q_kL(M^\star)\le F_R(q_k)$이므로
\[
\begin{aligned}
(q^\star-q_k)T(R)
&\le(q^\star-q_k)L(M^\star)\\
&=I(M^\star)-q_kL(M^\star)\\
&\le F_R(q_k).
\end{aligned}
\]
$T(R)>0$으로 나눈 뒤 양변에 $q_k$를 더하면
\[
q^\star
\le q_k+\frac{F_R(q_k)}{T(R)}.
\]
기존 상한 $u_k$와 위 식의 우변은 모두 $q^\star$ 이상이므로 상한을
\[
u_{k+1}
=\min\left\{
u_k,\;
q_k+\frac{F_R(q_k)}{T(R)}
\right\}
\]
로 갱신한다. 각 반복에서 $q_k\le q^\star\le u_k$를 유지하며, $q_0>0$일 때
\[
u_k-q_k\le\varepsilon q_0
\]
이면 반복을 종료한다.

구현에서는 float64 연산과 $\varepsilon=10^{-6}$을 사용한다. DP를 $K$회 호출할 때 DP 계산과 mask 복원에 드는 총 시간은 $O(K(N-R+1)(R+1))$이다.

\section{타일 기반 행 선택}\label{sec:tiles}

Tiles는 일정 길이의 타일을 단위로 행을 선택하고, 남은 예산을 하나의 연속 구간으로 채운다.

타일 길이를 양의 정수 $B$라 하자. 시작 행을 $0$과 $d=\lfloor B/2\rfloor$로 두면 두 배치의 타일은
\[
\boxed{
\begin{aligned}
\mathcal T_1&=\{C_{0,B},C_{B,B},C_{2B,B},\ldots\},\\
\mathcal T_2&=\{C_{d,B},C_{d+B,B},C_{d+2B,B},\ldots\}
\end{aligned}
}
\]
이다. 끝 행이 $N-1$을 넘는 타일은 제외하고, 각 배치에서 mask를 독립적으로 구성한다.

행 예산을
\[
R=KB+r,
\qquad
K=\left\lfloor\frac RB\right\rfloor,
\qquad
0\le r<B
\]
로 나누고, 각 배치에서 중요도 $\sum_{i\in C}v_i$가 큰 $K$개 타일을 선택한다.

$r>0$이면 선택한 타일과 겹치지 않는 길이 $r$의 연속 구간 중 중요도가 가장 큰 구간을 추가한다. $R<B$이면 중요도가 가장 큰 길이 $R$의 연속 구간을 선택한다.

예산을 채운 mask 중 $I(M)/L(M)$이 큰 것을 $M_{\mathrm{Tiles}}$로 선택한다.

$B\ge\ell_0$이고 $R\ge B$이면, two-line 모델의 추가 비용은 완전 타일과 합쳐지지 않은 잔여 구간에서만 발생할 수 있다. 따라서
\[
\alpha R
\le L(M_{\mathrm{Tiles}})
\le\alpha(R-r)+T(r).
\]
이다. $r=0$이거나 잔여 구간이 완전 타일과 합쳐지면 $L(M_{\mathrm{Tiles}})=\alpha R$이다. 한편 $B=1$이면 Tiles는 Top-$R$과 같다.

중요도 누적합과 상위 $K$개 타일의 부분 선택을 이용하기 때문에, 고정된 $B$에서 시간 복잡도는 평균 $O(N)$, 메모리 복잡도는 $O(N)$이다.

\section{최소 구간 길이 제약과 최적값의 보장}\label{sec:guarantee}

Two-line 모델에서는 같은 수의 행을 선택하고 모든 구간이 충분히 길면 읽기 비용이 동일해진다. 이 성질을 이용하여 길이 제약 아래에서의 최적화 문제를 정리하고, 중요도 보존율로 원래 최적값에 얼마나 가까운지 평가해 볼 것이다.

mask $M$의 중요도--지연시간 비율을 $q(M)=I(M)/L(M)$이라 하자.

\subsection{길이 제약 아래에서의 최적해}\label{sec:guarantee-constrained}

정확히 $R$개 행을 선택하는 문제에, 모든 최대 연속 구간 $C\in\mathcal C(M)$의 길이가 $\ell_0$ 이상이어야 한다는 조건을 추가하자. 이 조건이 가능하도록 $\lceil\ell_0\rceil\le R\le N$을 가정한다. 이때 길이 $R$의 단일 구간은 항상 조건을 만족한다.

각 구간은 $T(|C|)=\alpha|C|$를 만족하고, 구간 길이의 합은 $R$이므로 전체 읽기 비용은
\[
\begin{aligned}
L(M)
&=\sum_{C\in\mathcal C(M)}T(|C|)\\
&=\alpha\sum_{C\in\mathcal C(M)}|C|
=\alpha R.
\end{aligned}
\]
따라서 선택한 행의 위치나 구간의 수가 달라도 읽기 비용은 같다. 중요도--지연시간 비율의 분모가 고정되므로, 이 조건 아래에서는 중요도를 최대화하는 mask가 비율도 최대화한다.

그 mask를 $M_{\mathrm{sat}}^\star$, 보존 중요도를 $I_{\mathrm{sat}}=I(M_{\mathrm{sat}}^\star)$이라 하자. 길이 제약 아래에서의 최적 비율은
\[
q_{\mathrm{sat}}
=q(M_{\mathrm{sat}}^\star)
=\frac{I_{\mathrm{sat}}}{\alpha R}
\]
이다.

\subsection{일반적인 mask에 대한 보장}\label{sec:guarantee-general}

이제 같은 행 예산에서 길이 조건을 두지 않은 문제의 최적 비율 $q^\star$와 비교하자.

중요도가 가장 높은 $R$개 행을 선택한 Top-$R$ mask를 $M_T$, 그 중요도를 $I_T=I(M_T)$라 하자. 정확히 $R$개 행을 선택한 모든 mask는 $I(M)\le I_T$를 만족한다. 또한 two-line 모델에서 $L(M)\ge\alpha R$이므로
\[
q^\star\le\frac{I_T}{\alpha R}.
\]
이 상한은 중요도 상한 $I_T$와 비용 하한 $\alpha R$을 결합한 것이다. 두 값을 동시에 달성하는 mask가 존재하지 않더라도 성립한다.

$I_T>0$이면 $q^\star>0$이다. 정확히 $R$개 행을 선택하는 임의의 mask $M$에 대해 $q(M)\ge0$이므로
\[
\boxed{
\begin{aligned}
\frac{q(M)}{q^\star}
&\ge\frac{q(M)}{I_T/(\alpha R)}\\
&=\frac{I(M)}{I_T}\,
\frac{\alpha R}{L(M)}.
\end{aligned}
}
\]
오른쪽 마지막 식의 두 항은 각각 Top-$R$ 대비 중요도 보존율과 읽기 효율의 상한 대비 달성률이다. 후자는 행 단위 처리량 $R/L(M)$을 그 상한 $1/\alpha$로 나눈 값이다. 두 항은 모두 $0$과 $1$ 사이에 있으며, 그 곱으로 최적 비율 대비 달성률의 하한을 얻는다.

\subsection{중요도 보존율에 따른 보장}\label{sec:guarantee-importance}

모든 최대 연속 구간의 길이가 $\ell_0$ 이상인 mask는 비용 하한 $L(M)=\alpha R$을 달성한다. 따라서 앞 절의 보장에서 읽기 효율에 해당하는 항이 $1$이 되어
\[
\frac{q(M)}{q^\star}
\ge
\frac{I(M)}{I_T}
\underbrace{\frac{\alpha R}{L(M)}}_{=1}
=\frac{I(M)}{I_T}.
\]
이 조건 아래에서는 읽기 비용을 더 줄일 수 없으므로, 중요도 보존율만으로 최적 비율 대비 달성률을 보장할 수 있다. 중요도 상한 $I_T$에 도달하지 못하더라도, 비용 쪽의 하한은 이미 달성한 상태이기 때문이다.

특히 길이 제약 아래에서의 최적 mask $M_{\mathrm{sat}}^\star$에 적용하면
\[
\boxed{
\frac{q_{\mathrm{sat}}}{q^\star}
\ge
\frac{I_{\mathrm{sat}}}{I_T}.
}
\]
$M_{\mathrm{sat}}^\star$도 원래 문제에서 선택할 수 있는 mask이므로 $q_{\mathrm{sat}}\le q^\star$이다. 따라서 길이 제약 아래에서 Top-$R$ 중요도의 $95\%$ 이상을 보존하면, 중요도--지연시간 비율도 원래 최적값의 $95\%$ 이상이다.

이 보장은 원래 최적 mask를 직접 구하지 않고도 $I_{\mathrm{sat}}$과 $I_T$만으로 확인할 수 있다. 여기서 비용 하한의 달성과 비율에 대한 보장은 two-line 모델의 추정 비용을 기준으로 한다.
